% Level 3, Week 6 -- Review and Checkpoint
\lesson{6}{Review and Checkpoint}{Connect the pieces, and use the debugger to watch your code think, one line at a time.}
\lessoninfo{Concepts: Thonny debugger (Step over, Step into, Variables), OOP and algorithms (review)}{Tools: Thonny debugger, git}{CSTA: 3A-AP-21, 2-AP-17}

\begin{objectives}
  \item The student can use the Thonny debugger (Debug, Step over, Step into and the Variables view) to follow a program one line at a time. (L3.G8)
  \item The student can find and fix bugs in classes, sorts and searches using error messages, trace tables and the debugger. (L3.G8)
  \item The student can match each algorithm and data structure from weeks 1--5 to a problem it solves. (L3.G2)
  \item The student can describe a difficult bug, how they found it, and what they will try next time. (L3.G8)
\end{objectives}

\begin{materials}
  \item Print: the \textbf{Mid-level checkpoint} from the Assessments booklet, workbook chapter 6 and homework 6 (one per student).
  \item Match-up cards: copy the 10 pairs from the table below onto 20 index cards (or print this page and cut along the lines).
  \item On each Mac: Thonny with \texttt{3-advanced/exercises/week06} open. Turn on the Variables view: \menu{View > Variables} (a check mark appears), and also \menu{View > Object inspector} (it shows an object's attributes).
  \item Try it once yourself (3 minutes): open \texttt{ex-bug-missing-self/starter.py}, click the bug icon, then \menu{Step over} a few times.
\end{materials}

\section{At a glance}
\begin{timeplan}
  \timeslot{0--6}{Warm-up}{Unplugged \enquote{Algorithm match-up}: 10 names, 10 uses.}
  \timeslot{6--13}{Concept}{Tour of the Thonny debugger: Debug, Step over, Step into, Variables, Stop.}
  \timeslot{13--21}{Guided demo}{Debug Bug Hunt 1 (Score Keeper) together: three \enquote{forgot self} bugs.}
  \timeslot{21--52}{Hands-on build}{21--36: bug-hunt workshop, pairs work through Bug Hunts 2--4, swapping roles at each one. 36--51: \textbf{Mid-level checkpoint} (individual, quiet, 15 minutes). Early finishers: free-build.}
  \timeslot{52--60}{Wrap-up}{Growth-mindset reflection, exit ticket, homework.}
\end{timeplan}

\section{Warm-up (0--6 min)}
\begin{unplugged}[Algorithm match-up]
Shuffle the cards face up on the table. Students match each name with its use, then explain one match each. With 1 student, time them and try to beat the time at the end of the session.
\begin{center}
\small\renewcommand{\arraystretch}{1.2}
\begin{tabular}{|p{2.8cm}|p{3.9cm}|p{2.8cm}|p{3.9cm}|}
\hline
\textbf{Name} & \textbf{Use} & \textbf{Name} & \textbf{Use} \\ \hline
git commit & a save point with a message & binary search & open a sorted dictionary in the middle \\ \hline
class & a blueprint for many objects & bubble sort & neighbors swap until big items reach the end \\ \hline
inheritance & a Sprinter is an Enemy & sorted(key=...) & order a leaderboard by points \\ \hline
composition & a Game has a list of Players & stack & the undo button: last in, first out \\ \hline
linear search & check every locker, one by one & queue & taking turns: first in, first out \\ \hline
\end{tabular}
\end{center}
\end{unplugged}

\section{Concept (6--13 min)}
\begin{teacherscript}
\enquote{When code goes wrong, it is usually because the computer is doing exactly what we wrote, not what we meant. A \textbf{debugger} lets us slow the computer down and watch it think, one line at a time, while we see every variable change. Professional programmers use one every day.}
\end{teacherscript}
Show these on your Mac while the students watch (Thonny's toolbar is at the top):
\begin{itemize}[itemsep=2pt]
  \item \textbf{Debug} (the bug icon, next to the green Run button): starts the program in slow motion. The first line is highlighted.
  \item \textbf{Step over} (or \keys{F6}): run the highlighted line and stop at the next one. Function calls run in one go.
  \item \textbf{Step into} (or \keys{F7}): go \emph{inside} a function or method call to watch it line by line. Thonny opens a small window for the function: that is normal, it shows the function's own variables.
  \item \textbf{Resume} (or \keys{F8}): run normally until the end (or the next breakpoint). \textbf{Stop} (the red stop sign): end the program.
  \item \textbf{Variables view}: every variable and its value right now. Click an object there and the \textbf{Object inspector} shows its attributes.
  \item \textbf{Breakpoint}: click on a line number to put a red dot there. Debug, then Resume, and Thonny pauses at that line.
\end{itemize}
On some Macs the F-keys need \keys{fn} held down, so clicking the toolbar buttons is easiest.

\begin{analogy}[A slow-motion replay]
Running a program is like watching a goal in a football match: it is over in a second. The debugger is the slow-motion replay with the score shown on screen: you can pause on every frame, see exactly where every player (variable) is, and spot the moment it went wrong.
\end{analogy}

\section{Guided demo (13--21 min)}
\begin{workedexample}[Debugging Bug Hunt 1: the Score Keeper]
Open \texttt{week06/ex-bug-missing-self/starter.py}. Students tell you what to do at each step.
\begin{enumerate}[itemsep=2pt]
  \item Press \textbf{Run}. Read the last line of the red error together: \texttt{TypeError: ScoreKeeper.add() takes 1 positional argument but 2 were given}. Say: \enquote{Python passed the keeper itself \emph{and} 10, but \texttt{add} only has room for one.} Fix: \texttt{def add(self, points):}.
  \item Press \textbf{Debug}. Click \textbf{Step over} until the line \texttt{keeper.undo\_last()} is highlighted. Look at the Variables view: click \texttt{keeper}, and the Object inspector shows \texttt{score} = 15 and \texttt{history} = [10, 5].
  \item Click \textbf{Step into} to go inside \texttt{undo\_last}. Step over until it crashes on \texttt{score -= last}: \texttt{UnboundLocalError}. Ask: \enquote{Whose score?} Fix: \texttt{self.score -= last}.
  \item Run again: \texttt{NameError: name 'nickname' is not defined} in \texttt{report}. Fix: \texttt{self.nickname}.
  \item Run once more: \texttt{Comet has 10 points}. Then open \texttt{check\_my\_work.py} and press Run: all 4 checks pass.
\end{enumerate}
\end{workedexample}

\section{Hands-on build (21--52 min)}
\begin{pairroles}[Detective and debugger driver]
The \textbf{detective} reads the error message aloud and makes a \emph{prediction}: \enquote{I think the bug is on line \ldots because \ldots}. The \textbf{driver} uses the debugger to test the prediction. Swap roles at each new bug hunt.
With \textbf{1 student}: you are a \enquote{rubber duck}: they explain each line to you before stepping over it. With \textbf{3}: the third student keeps a \textbf{bug log} on paper (bug, line, fix) and rotates in. With \textbf{4}: two pairs, starting at different bug hunts.
\end{pairroles}

\begin{bughunt}[Bug-hunt workshop (21--36 min)]
Each pair works through these folders in \texttt{week06}. Every starter is broken on purpose; it passes \texttt{check\_my\_work.py} only when fixed. About 5 minutes each.
\begin{description}[itemsep=2pt,font=\normalfont\bfseries]
  \item[Bug Hunt 2, Lost Scores] (\texttt{ex-bug-lost-items}): a sort that loses items. Hint if stuck: watch \texttt{scores}, \texttt{remaining} and \texttt{result} in the Variables view. Three bugs: no copy of the list, the loop stops one too early, duplicates skipped.
  \item[Bug Hunt 3, Word Finder] (\texttt{ex-bug-binary-search}): an off-by-one binary search. Hint: a trace table for \texttt{find\_word(["ant", "bee", "cat"], "ant")}. Two bugs: \texttt{high = len(words)} and \texttt{while low < high}. Fixing only one may cause an \texttt{IndexError}: that is progress!
  \item[Bug Hunt 4, Shape Party] (\texttt{ex-bug-shape-party}): inheritance bugs. First nothing is drawn (\texttt{shape.draw} without parentheses), then a crash (no \texttt{super().\_\_init\_\_}), then diamonds that are dots (\texttt{drow} typo).
\end{description}
Rule for the workshop: change \textbf{one} thing, run again. If two things change at once, you will not know which one fixed it.
\end{bughunt}

\textbf{Checkpoint (36--51 min).} Laptops closed; each student does the Mid-level checkpoint (Assessments booklet) alone and quietly. Read its instructions aloud, then say: \enquote{This shows me what to practice next; \enquote{not sure} is a fine answer.} If it runs over 15 minutes, finish it next week. Score it later with the answer key.

\begin{freebuild}[For early finishers]
When a student hands in the checkpoint early, they open Bug Hunt 4 (Shape Party) and remix the fixed version: new shape classes, a different arrangement, their own colors from the palette. Or they write a tiny buggy program of their own for a partner to debug next week.
\end{freebuild}

\section{Wrap-up (52--60 min)}
\begin{safety}[Growth mindset: bugs are how we learn]
Ask each student to share \enquote{my hardest bug today, and what helped}. Celebrate effort and strategy, not speed: \enquote{You changed one thing at a time: that is what professionals do.} Remind them that every programmer meets bugs every day, and that asking for help with a clear description (what I expected, what happened, what I tried) is a skill, not a weakness. Being kind when a partner's code is broken is part of being a good teammate.
\end{safety}

\begin{formativecheck}[Exit ticket]
\textbf{Look for:} students reading the last line of an error first; using Step into on method calls; checking the Variables view rather than guessing.\\
\textbf{Ask:} \enquote{What is the difference between Step over and Step into?}\\
\textbf{Answer:} Step over runs the whole line (including any function it calls) in one go; Step into goes inside the function so you can watch its lines one by one.
\end{formativecheck}

\begin{reflection}[Write on the workbook page]
\begin{itemize}[itemsep=1pt]
  \item Which strategy helped most today: the error message, a trace table, the debugger, or explaining to someone?
  \item Which topic from weeks 1--5 do you feel strongest in? Which one do you want to practice?
\end{itemize}
\end{reflection}

\section{Notes for the teacher}
\begin{commonmistakes}
  \item \textbf{Pressing Run instead of Debug}: the program runs at full speed. The debugger only starts with the bug icon.
  \item \textbf{Stepping over a method call} and missing the bug inside it: use Step into on lines that call your own methods.
  \item \textbf{Changing many lines at once} after an error. Fix one thing and run again.
  \item \textbf{Reading only the first line of a traceback}: the most useful line is the \emph{last} one, and the line number right above it.
  \item \textbf{Lost in the step-into windows, or no Variables view}: press Stop and Debug again; turn on \menu{View > Variables}.
\end{commonmistakes}

\begin{differentiation}
\supportpath Do the bug hunts as a whole group on one screen, with you as the driver and the students as detectives. Give the hints from the bug-hunt box straight away. Accept partial fixes.
\extensionpath After the workshop, set a breakpoint in \texttt{find\_word} and use only Resume to jump from loop to loop. Then write a new buggy function for a partner (one bug only) and swap.
\end{differentiation}

\begin{stretchchallenge}[Write a test that catches the bug]
Pick one bug from today. Write your own extra test function in that folder's \texttt{test\_*.py} file that fails with the buggy code and passes with the fixed code, for example \texttt{assert sort\_scores([2, 2]) == [2, 2]}. Run \texttt{check\_my\_work.py} to see it.
\end{stretchchallenge}

\begin{teachernote}[If you are short on time]
The checkpoint is the priority: keep its 15 minutes. Shorten the match-up to 5 pairs and the workshop to Bug Hunts 2 and 3. Bug Hunt 4 becomes a free choice for next week.
\end{teachernote}
